% Originally: content/week-04.tex, \section{The differential} (Corsin Week 4, Monday)

% Source: Corsin Nick/Class Notes/Week 4.pdf, p. 2
\section{The differential}
\label{sec:the_differential}

For $f : \mathbb{R}\to\mathbb{R}$ differentiable at $x_0 \in \mathbb{R}$:
\[f(x_0+h) = f(x_0) + f'(x_0)h + o(h)\]
\[\Updownarrow\]
\[\lim_{h\to 0}\frac{f(x_0+h)-f(x_0)}{h} = f'(x_0)\]
\[\Updownarrow\]
\[\lim_{h\to 0}\left|\frac{f(x_0+h)-f(x_0)-f'(x_0)h}{h}\right| = 0\]

For $F : \mathbb{R}^n \to \mathbb{R}^m$ differentiable at $x_0 \in \mathbb{R}^n$:
\[F(x_0+h) = F(x_0) + \underbrace{DF_{x_0}(h)}_{\text{linear map}} + o(|h|)\]
\[\Updownarrow\]
\[\lim_{h\to 0}\frac{|F(x_0+h)-F(x_0)-DF_{x_0}(h)|}{|h|} = 0\]
\[\Downarrow\]
\[\lim_{s\to 0}\frac{F(x_0+sv)-F(x_0)}{s} = DF_{x_0}(v) = \partial_vF(x_0) \quad \text{for each fixed } v \in \mathbb{R}^n\]
\[\text{(the \newterm{directional derivative} of } F \text{ at } x_0 \text{ in direction } v\text{)}\]

Note that the last arrow points \textbf{one way only}. The first two lines say the same thing;
the third is a strictly weaker consequence, and recovering differentiability from it is exactly
what fails in several variables --- see \cref{ex:partial_derivatives_not_continuous} below.

% Transition: Claude Opus 5 (High)
It is worth pausing on why the definition had to be restated rather than merely copied. The
one-variable derivative is defined as a \emph{quotient}, and that is precisely what does not
survive: for $h \in \mathbb{R}^n$ with $n \geq 2$ the expression
$\big(F(x_0+h)-F(x_0)\big)/h$ is meaningless, because one cannot divide by a vector. The first
line above is the version that does survive, and it is the one to remember: \textbf{$F$ is
differentiable at $x_0$ if it is approximated near $x_0$ by an affine map, to an error smaller
than $|h|$.} Everything in this chapter --- the Jacobian, the chain rule, Taylor expansions ---
is bookkeeping for that single idea, and the reason the derivative becomes a \emph{linear map}
rather than a number is simply that a linear map is what an approximation in $n$ variables looks
like.

\begin{ainote}
Corsin's middle line reads
$\lim_{h\to0}\frac{F(x_0+h)-F(x_0)}{|h|} = DF_{x_0}\!\left(\frac{h}{|h|}\right)$, inside a chain
of $\Updownarrow$. As written that limit does not exist: $h$ is being sent to $0$ on the left
while still appearing on the right, and the quotient genuinely depends on the direction of
approach. The intended statement is the one typeset above --- fix a direction $v$ first, then
let the scalar $s \to 0$ --- and it is only an \emph{implication}, not an equivalence. The
distinction is not pedantic: it is the entire content of
\Cref{ex:partial_derivatives_not_continuous} and of
\Cref{rem:directional_derivative_caveat}.
\end{ainote}

% Source: Corsin Nick/Class Notes/Week 4.pdf, p. 2
% (Re-cited: the paragraph and AI-Note above are additions, not transcription.)
\begin{definition}[Differentiability]
\label{def:differentiable}
We say that $F : \mathbb{R}^n \to \mathbb{R}^m$ is \newterm{differentiable}
(\germanterm{differenzierbar}) at $x_0$ if a linear map $DF_{x_0}$ as above exists, where
$x_0, h \in \mathbb{R}^n$. The same definition applies verbatim to $F : U \to \mathbb{R}^m$
for any open subset $U \subseteq \mathbb{R}^n$, and we use it in that generality throughout.
\end{definition}

% Sonnet 5 (Medium)
\begin{remark}
The linear map $DF_{x_0}$ in \cref{def:differentiable}, when it exists, is automatically unique:
if $L_1, L_2$ both satisfy the defining limit, then
$(L_1-L_2)(h)/|h| \to 0$ as $h \to 0$ along every direction, which forces the linear map
$L_1 - L_2$ to vanish identically. Concretely: write $L := L_1-L_2$ and suppose $L(h) \neq 0$
for some $h$. Along the ray $th$, $t > 0$, linearity gives
\[\frac{|L(th)|}{|th|} = \frac{t\,|L(h)|}{t\,|h|} = \frac{|L(h)|}{|h|} \neq 0,\]
a nonzero \emph{constant} --- so the quotient cannot tend to $0$ as $t \to 0$. Hence $L(h)=0$
for every $h$. This
is the multivariable analogue of the uniqueness-of-limits argument from
\Cref{prop:uniqueness_of_limits}.
\end{remark}

\begin{proposition}[Differentiability implies continuity]
\label{prop:differentiable_implies_continuous}
If $F : U \to \mathbb{R}^m$, $U \subseteq \mathbb{R}^n$ open, is differentiable
(\cref{def:differentiable}) at $x_0 \in U$, then $F$ is continuous at $x_0$.
\end{proposition}

\begin{proof}
By \cref{def:differentiable}, $F(x_0+h) - F(x_0) = DF_{x_0}(h) + o(|h|)$ as $h \to 0$. Since
$DF_{x_0}$ is linear on the finite-di\-men\-sional space $\mathbb{R}^n$, it is bounded, i.e.\ there
exists $C \geq 0$ with $|DF_{x_0}(h)| \leq C|h|$ for all $h$. Hence
\[|F(x_0+h) - F(x_0)| \leq |DF_{x_0}(h)| + o(|h|) \leq C|h| + o(|h|) \xrightarrow{h\to 0} 0,\]
so $F(x_0+h) \to F(x_0)$ as $h \to 0$, i.e.\ $F$ is continuous at $x_0$.
\end{proof}

\begin{remark}
The converse fails: continuity does not imply differentiability, exactly as in one variable
(e.g.\ $x \mapsto |x|$ at $0$). \Cref{ex:partial_derivatives_not_continuous} below shows the
gap is even wider in several variables --- \emph{all} partial derivatives can exist at a point
without $F$ even being continuous there, let alone differentiable.
\end{remark}
